% Ad Atticum 16.11 (ad-atticum-16-11) --- English translation
% Translated by Claude (Anthropic) from the Latin (Perseus).
% Style guide: STYLE.md.

\ciceroLetterOpener{CICERO ATTICO salutem}

\ciceroSection{1} On the Nones I received two letters from you, one that you had dispatched on the Kalends, the other on the day before. So, first to the earlier. I am glad that you approve of my work; from it you culled the flowers \textit{[Greek: anthē]} themselves --- they looked the more brilliant to me for your judgement, since I had been dreading those little vermilion-wax marks of yours. About Sicca, it is as you write; I scarcely held myself back from a sharp word. So I will graze him with no insult to Sicca or Septimia, only enough that they may know he had children's children \textit{[Greek: paides paidōn]} from the daughter of Gaius Fadius --- without the Lucilian phallus \textit{[Greek: phallos]}. And how I wish I might see the day when that speech shall roam at large so freely that it makes its way even into Sicca's house! --- but we should need the times the Triumvirs had for that. May I die if it isn't witty! You, though, will read it to Sextus and write me back his judgement in full. One man, to me, is ten thousand \textit{[Greek: heis emoi murioi]}. You will guard against the intrusion of Calenus and of Calvena.

\ciceroSection{2} As for your fear of being a chatterbox \textit{[Greek: adoleschos]} with me --- who is less so? To me, as the iambic of Archilochus did to Aristophanes, so each of your letters seems best the longer it is. As for your putting me right --- why, even if you were rebuking me I would not merely bear it gladly but rejoice in it, since in your rebuke there is judgement together with good will \textit{[Greek: eumeneia]}. So I shall gladly correct those things you have marked, by the same right as in the case of Rubrius rather than that of Scipio; and I shall pare away the heaping praise of Dolabella. And yet there is in that passage a pretty irony \textit{[Greek: eirōneia]}, it seems to me, where I have him three times in battle-line against fellow citizens. I prefer also this: ``Nothing is more disgraceful than that this man should live'' --- or what could be more disgraceful?

\ciceroSection{3} I am not put out that Varro's \textit{Peplography} \textit{[Greek: Peplographia]} meets with your approval; from him I have not yet got that \textit{Heraclidean} piece \textit{[Greek: Hērakleideion]}. You urge me to write --- as a friend, indeed; but know that I am doing nothing else. Your head-cold troubles me. Please, apply the care you usually do. ``O Titus''\,---\,I am glad I do you good. The Anagnines are Mustela the squadron-leader \textit{[Greek: taxiarchēs]} and Laco who drinks the most. The book you ask for I will polish off and send.

\ciceroSection{4} This for the later letter. The topic of \textit{[Greek: ta peri tou]} how far Panaetius went --- I have finished it in two books. He has three; but although at the outset he had laid down the division thus --- that there are three classes of inquiry concerning duty: one, when we deliberate whether a thing is honourable or base; a second, whether it is useful or useless; a third, when these seem to fight with one another, how judgement is to be made (Regulus's case is of this kind: to return was honourable, to stay was useful) --- on the first two he discoursed splendidly; on the third he promises that he will follow in due course, but he wrote nothing. That topic Posidonius pursued. I have both sent for his book and written to Athenodorus Calvus to send me the headings \textit{[Greek: ta kephalaia]}; I am waiting for them. Please urge him on and ask him to send them as soon as possible. In that work is the section on circumstance \textit{[Greek: peri tou kata peristasin]}. As for the title you ask about --- I have no doubt that \textit{kathēkon} \textit{[Greek: kathēkon]} is ``officium,'' unless you think otherwise; but a fuller title is \textit{De Officiis}. I am dedicating \textit{[Greek: prosphōnō]} the work to my son Cicero. It seemed not inappropriate \textit{[Greek: anoikeion]}.

\ciceroSection{5} About Myrtilus, you write plainly. Oh, what a man you always are with people like that! Really? --- against Decimus Brutus? May the gods give them what they deserve!

\ciceroSection{6} I, as I had written, did not bury myself in the Pompeianum --- first because of the storms, than which nothing was more foul; then because of daily letters from Octavian urging me to take up the business, to come to Capua, to save the republic a second time, and in any case to come to Rome at once: ``\textit{they were ashamed to refuse}'' \textit{[Greek: aidesthen men anēnasthai]}. Still, he is acting with real energy, and goes on acting. He will come to Rome with a great body of men, but he is plainly a boy. He supposes the senate will meet at once. Who will come? If anyone does, who will give offence to Antony in such uncertain times? On the Kalends of January he may perhaps be of use --- or rather there will be a pitched battle before that. The country towns are wonderfully partial to the boy. On his march into Samnium he came to Cales and stayed at Teanum: a remarkable welcome \textit{[Greek: apantēsis]} and call to arms. Would you have thought it? On account of this I shall be at Rome sooner than I had decided. As soon as I have decided, I will write.

\ciceroSection{7} Although I have not yet read the bond-terms --- for Eros had not arrived --- still, I should be glad if you could settle the matter on the day before the Ides. Letters to Catina, Tauromenium, and Syracuse I shall be able to send more conveniently if the interpreter Valerius will write me the names of men of influence there. The same men are not in favour at all times, and our own associates are mostly dead. Still, I have written officially, in case Valerius wished to use them; or else let him send me names.

\ciceroSection{8} About the Lepidian holidays Balbus has written to me. I shall expect your letter up to the third day before the Kalends and reckon I shall then know about Torquatus's little business. Quintus's letter I have sent on to you so that you might see how much he loves the man whom he is sorry you love less. As for Attica, since she is in the bright spirits that are the best thing in children, please give her a kiss from me.
